\documentclass{beamer}

% +==============================+
% | Vydefinování maker a packagů |
% +==============================+

\usepackage[utf8]{inputenc}
\usepackage[czech]{babel}
\usepackage{graphicx}
\usepackage{tikz}
\usetikzlibrary{shapes.arrows}
\usepackage{MnSymbol,wasysym}
\usepackage{multirow}
\usepackage{xcolor}

\newcommand{\cevroLogo}{
	\begin{tikzpicture}[remember picture,overlay]  
		\node [xshift=-1.6cm,yshift=-0.5cm] at (current page.north east)
		{\includegraphics[width=3cm,height=2cm,keepaspectratio]{images/newcevro.png}};
	\end{tikzpicture}
}

\newcommand{\specialcell}[2][c]{%
  \begin{tabular}[#1]{l@{}c@{}}#2\end{tabular}}
  
\newcommand\Wider[2][3em]{%
\makebox[\linewidth][c]{%
  \begin{minipage}{\dimexpr\textwidth+#1\relax}
  \raggedright#2
  \end{minipage}%
  }%
}

% +=========================+
% | Vydefinování tikz stylů |
% +=========================+

\tikzset{
    myarrow/.style={
        draw,
        fill=green,
        single arrow,
        minimum height=3.5ex,
        single arrow head extend=1ex
    }
}

\tikzset{
    myarrowred/.style={
        draw,
        fill=red,
        single arrow,
        minimum height=3.5ex,
        single arrow head extend=1ex
    }
}

% +=======================================+
% | Základní informace o bakalářské práci |
% +=======================================+

\title{Svěřenský fond jako nástroj mezigeneračního převodu majetku}
%\subtitle[krátký název prezentace]{dlouhý název prezentace}
%\author{Dominik Bálint}
\author[My name]{Dominik Bálint\\ dominik.balint@vsci.cz \\ 775~063~885 \\ \hspace*{1mm} \\\footnotesize Vedoucí práce: Mgr. Ing. Tomáš Střeleček, LL.M., Ph.D.\\ \footnotesize Oponent práce: JUDr. Miroslav Sedláček, Ph.D., LL.M.}
\institute{CEVRO Institut}
\date{\today}

% +=============+
% | První sekce |
% +=============+

\section{Prezentace k obhajobě}

% +===============================================+
% | Použitý theme v prezentaci k bakalářské práci |
% +===============================================+

\usetheme{Singapore}

\begin{document}

% TO DO:
% 	Slide - kontaktní údaje 				- [x]
% 	Slide - název práce 					- [x]
% 	Slide - klíčová slova v ČJ a ENG 	- [x]
% 	Strukturu práce: -|
% 		Předmluva 						- [ ]
% 		Úvod 							- [x]
% 		Zdroje 							- [x]
% 		Vlastní text práce 				- [x]
% 		Závěr 							- [x]
% 		Resumé 							- [x]
% 		Přílohy 						- [x]
% 	A hlavně: -|
% 		Anotaci 						- [x]
% 		a Abstrakt 						- [x]

% +===============+
% | Titulní slide |
% +===============+

\maketitle

% +====================+
% | Začátek prezentace |
% +====================+

\begin{frame}
	\frametitle{Obsah}
	\cevroLogo

\begin{itemize}
	\item Anotace
	\item Abstrakt
	\begin{itemize}
		\item Cíl práce%\pause
		\item Klíčová slova%\pause
		\item Metodika%\pause
		\item Struktura práce
		\item Výsledky, doporučení%\pause
		\item Závěr%\pause
	\end{itemize}
	\item Zodpovězení otázek
\end{itemize}
	
\end{frame}

% +=============+
% | První slide |
% +=============+

\begin{frame}
	\frametitle{Anotace}
	\cevroLogo
	
Práce je zaměřena na problematiku svěřenských fondů. Zabývá se historickými souvislostmi vedoucími ke společenské potřebě vytvořit instituty podobné svěřenskému fondu a následným vývojem těchto institutů završených přijetím právní úpravy quebecké fiducie do českého právního řádu ve formě svěřenského fondu.\\

\vspace*{5mm}

Práce dále klade důraz na definování tohoto institutu, jeho právní kodifikaci a charakteristiku ve vztahu k jeho užití pro mezigenerační převod majetku.
		
\end{frame}


% +=============+
% | Druhý slide |
% +=============+

\begin{frame}
	\frametitle{Cíl práce}
	\cevroLogo

\fbox{Cílem práce} jest představit svěřenský fond,\\
jeho použití k mezigeneračnímu převodu majetku,
a zodpovězení \textbf{hypotézy}, zda lze svěřenský fond úspěšně použít k obcházení práva nepominutelných dědiců na povinný díl.

\vspace*{5mm}

\fbox{Vedlejší cíle} : další problémy spojené se svěřenským fondem.

\end{frame}

% +=============+
% | Třetí slide |
% +=============+

\begin{frame}
	\frametitle{Klíčová slova}
	\cevroLogo	
	
\fbox{Klíčová slova} : {\footnotesize Trust, svěřenský fond, svěřenství, \textit{inter vivos}, \textit{mortis causa}, dědictví, dědic, nepominutelný dědic, smlouva, pořízení pro případ smrti, vlastnictví, fideikomis, Řím, římské právo, kontinentální právo, angloamerické právo, odkaz, Česká republika, Anglie, Kanada, Lichtenštejnsko, Německo, Itálie, Francie, Quebec, CCQ, občanský zákoník, důkaz, zakladatelské právní jednání, statut, vznik, novelizace, právní subjektivita, právnická osoba, bezsubjektové vlastnictví, mezigenerační převod majetku, darování, dědická sukcese, daňové zatížení, věřitelé.}

\vspace*{1mm}

\fbox{Key words} : {\footnotesize fiducie, trusteehip, inheritance, heir, forced heirship, contract, last will and testament, ownership, Rome, roman law, continental law, common law, legacy, Czech republic, England, Canada, Lichtenstein, Germany, Italy, France, civil code, evidence, instruments, statute, founding, amendment, legal subjectivity, legal entity, non-subject ownership, intergenerational property transfer, gift, hereditary succession, tax burden, creditors.}

\end{frame}

% +==============+
% | Čtvrtý slide |
% +==============+

\begin{frame}
	\frametitle{Zdroje}
	\cevroLogo
	
\begin{enumerate}
\item BEDNAŘÍKOVÁ Barbora. \textit{Svěřenské fondy. Institut pro uchování a převody rodinného majetku.} 2., aktualizované vyd. Praha: Wolters Kluwer, 2014.
\item DVOŘÁK J. - ŠVESTKA J. - ZUKLÍNOVÁ M. a kol. \textit{Občanské právo hmotné.} Svazek 1. Díl: Obecná část. 2.,aktualizované a doplněné vydání. Praha: Wolters Kluwer, 2016.	
\item ELIÁŠ Karel a kolektiv. \textit{Nový občanský zákoník s aktualizovanou důvodovou správou a rejstříkem.} Nakladatelství Sagit, a.s., 2012. isbn: 978-80-7208-922-2.
\item HAYTON D. J. \textit{Extending the Boundries of Trusts and Similar Ring-Fenced Funds.} Kluwer Law International, 2002. isbn: 978-90-411-9879-2.
\item PIHERA V. \textit{Nejpodivnější zvíře v lese – poznámky ke svěřenskému fondu.} Vol. č. 10. Obchodněprávní revue., 2012.
\end{enumerate}

\end{frame}

% +==============+
% | Čtvrtý slide |
% +==============+

\begin{frame}
	\frametitle{Metodika}
	\cevroLogo

\begin{columns}[T]
	\column{0.5\textwidth}
	\textbf{Stanovení cíle}\\
		Z pozitivistického hlediska prozkoumat vztah mezi jednotlivými ustanoveními svěřenského fondu a dědického práva.\\
		---\\
		Účelem dovození odpovědi na vyloženou hypotézu a odhalení problémů.\\
		---\\
		Konstruování normativního výstupu - navržení obecných změn, pokud jich bude třeba.
	\column{0.5\textwidth}
	\textbf{Výzkumné metody}\\
		Obecně teoretické metody.\\
		---\\
		Popisná, kompilační metoda u teoretické části.\\
		---\\
		Analyticko-syntetické metody u praktické části.\\
		---\\
		Zkoumáno (zkoumaná problematika) dvěma optikami ve smyslu interpretace právních norem.\\
		%\vspace*{5.9cm}
\end{columns}

\end{frame}

% +==============+
% | Pátý slide   |
% +==============+

\vspace*{-1cm}

\begin{frame}
	\frametitle{Metodika}
	\cevroLogo
	
\includegraphics[height=\textheight,width=\textwidth,keepaspectratio]{images/Zkoumani.png}

\vspace*{0.2cm}

\begin{columns}[T]
\column{0.5\textwidth}
{\small\textbf{Optika 1}}\\
Systematická\\
Logická\\
Teleologická\\
Komparativní (pro ilustraci vzorových řešení)\\
Doktrinální\\
Smysl ve vztahu ke konstrukci (struktuře) svěřenského fondu.\\
\column{0.5\textwidth}
{\small\textbf{Optika 2}	}\\
Formální - povrchní\\
Gramatická, §2\\
"Logická"\\
Dedukce úmyslu zákonodárce\\
\end{columns}

	
\end{frame}

% +===============+
% | Šestý slide   |
% +===============+

\begin{frame}
	\frametitle{Struktura práce}
	\cevroLogo

\Wider{
\begin{table}[]
\begin{tabular}{|l|l||l||l||l||l|}
	\hline
			\textbf{Sekce} & 
			\textbf{\specialcell[t]{1. Resumé a\\klíčová slova}} & 
			$\rightarrow$ & 
			\textbf{\specialcell[t]{2. Úvod}} & 
			$\rightarrow$ & 
			\textbf{\specialcell[t]{3. Vlastní\\text práce}} \\
		\hline
			\textbf{Podsekce} & 
			\specialcell[t]{Resumé\\Summary\\Klíčová slova} & 
			&
			\specialcell[t]{Úvod} & 
			&
			\specialcell[t]{Obecný úvod\\Teoretická část\\Praktická část} \\
	\hline
	\hline
			\textbf{Sekce} & 
			\textbf{\specialcell[t]{4. Závěr}} & 
			$\rightarrow$ & 
			\textbf{\specialcell[t]{5. Zdroje}} & 
			$\rightarrow$ & 
			\textbf{\specialcell[t]{6. Přílohy}} \\
		\hline
			\textbf{Podsekce} & 
			\specialcell[t]{Závěr} & 
			&
			\specialcell[t]{Zdroje} & 
			&
			\specialcell[t]{Vzor smlouvy\\Vzor listiny} \\
	\hline
\end{tabular}
\end{table}}
	
\end{frame}

% +=============+
% | Sedmý slide |
% +=============+

\begin{frame}
	\frametitle{Výsledky, doporučení}
	\cevroLogo

\fbox{Cíle jednotlivých částí}\\
\vspace*{5mm}
Teoretická část $\rightarrow$ popisná část (historie, přijetí, právní úprava).\\
Praktická část $\rightarrow$ analytická část (§1400 - §1474 OZ {\color{red} X} §1475 - §1720 OZ).\\
Závěrečná část $\rightarrow$ návrh řešení problematických částí (pokud budou nalezeny - především kolize s dědickým právem).\\
\vspace*{5mm}
\fbox{Výsledek analytické částí}\\
\vspace*{5mm}
Ve vztahu k hypotéze {\color{red} nedosaženo} výsledku $\rightarrow$ hypotéza nezodpovězena, vysloven pouze osobní názor, domněnka ochrany nepominutelných dědiců $\rightarrow$ dovození problematiky, potvrzení problémů, shoda {\color{red} X} konfrontace s literaturou zaměřující se na stejné téma.\\

\end{frame}

% +=============+
% | Osmý slide  |
% +=============+

\begin{frame}
	\frametitle{Výsledky, doporučení}
	\cevroLogo

\textbf{Výkladová optika} $\rightarrow$ {\color{red} nedovozená žádná} ochrana.\\
\textbf{Úmyslová optika} $\rightarrow$ {\color{green} potenciálně dovozená} ochrana, kolize s absencí právní osobnosti svěřenského fondu popřípadě její fikce ve vztahu k pozůstalostnímu řízení, {\color{red} !} nejedná se o dovození ochrany, pouze o dovození záměru zákonodárce ochranu vtělit do daného právního předpisu.\\
-----------------------------------------\\
{\color{red} Proti}: absence právní osobnosti svěřenského fondu, svěřenskému správci nesvědčí vlastnické právo, vyčleněný majetek není součástí pozůstalosti, autonomie vůle, potenciální znevýhodnění svěřenského fondu \textit{mortis causa} oproti svěřenskému fondu \textit{inter vivos}.\\
{\color{green} Pro}: §1452, §311, §1492 OZ, částečně také §17, odstavec 2.\\
Doporučení: novelizace.

\end{frame}

% +==============+
% | Devátý slide |
% +==============+

\begin{frame}
	\frametitle{Závěr}
	\cevroLogo

\setbeamercolor{block title}{use=structure,fg=white,bg=blue!75!black}
\setbeamercolor{block body}{use=structure,fg=black,bg=white!20!white}
\vspace*{-5mm}
\begin{block}{\small§ 1452}\small{
(1) Každý svěřenský fond musí mít statut. Statut svěřenského\\ fondu vydává zakladatel. Zřizuje-li se svěřenský fond\\ pořízením pro případ smrti, použije se přiměřeně § 311.}
\end{block}
\vspace*{-2mm}
\centering\tikz [baseline=-1ex]{\node [myarrow,rotate=-90] {};}\\
\vspace*{-2mm}
\begin{block}{\small§ 311}\small{
(1) Při založení nadace pořízením pro případ smrti se do nadace vnáší vklad povoláním nadace za dědice nebo nařízením odkazu. V takovém případě nabývá založení nadace účinnosti smrtí zůstavitele.}
\end{block}
\vspace*{-2mm}
\centering\textbf{\Large{\color{green} +}}\\
\vspace*{-2mm}
\begin{block}{\small§ 1492}\small{
Pořízením pro případ smrti nelze zkrátit povinný díl nepominutelného dědice, který se práva na povinný díl nezřekl a nedošlo-li ani k vydědění. Pokud tomu pořízení pro případ smrti odporuje, náleží nepominutelnému dědici povinný díl.}
\end{block}

\end{frame}

% +==============+
% | Desátý slide |
% +==============+

\begin{frame}
	\frametitle{Závěr}
	\cevroLogo

\begin{center}
Absence kolizní právní úpravy,\\
doktrinálně odlišné pohledy,\\
absence soudních rozhodnutí.\\
\tikz [baseline=-1ex]{\node [myarrowred,rotate=-90] {};}\\
Právní nejistota,\\
není jasné, zda je dědic chráněn,\\
nutnost explicitní úpravy.
\end{center}

\end{frame}


% +=================+
% | Jedenáctý slide |
% +=================+

\begin{frame}
	\frametitle{Zodpovězení otázek}
	\cevroLogo

\begin{table}[]
\begin{tabular}{|l|}
\hline
Otázka 1 Vedoucí práce                                                                                                                                                                                                                       \\ \hline
\begin{tabular}[c]{@{}l@{}}Které otázky problematického vztahu právní úpravy svěřenských\\ fondů a dědického práva si podle autora žádají nejnaléhavěji\\ řešení přijetím novely?\\ Jak by taková novela mohla, podle autora, znít?\end{tabular} \\ \hline
Nejzásadnější změna by měla nastat ve smyslu kolize dědického\\ práva s právem svěřenských fondů.\\
---\\
Další problémy vytyčené v práci (započtení, daňová problematika)\\ jsou oproti této problematice absolutně marginálního charakteru.\\
---\\
Mělo by dojít primárně k řešení problematiky nepominutelných\\ dědiců.                                                                                                                                                                                                                             \\ \hline
\end{tabular}
\end{table}

\end{frame}

% +================+
% | Dvanáctý slide |
% +================+

\begin{frame}
	\frametitle{Zodpovězení otázek}
	\cevroLogo

Pro účely vyřešení kolizní problematiky stačí považovat svěřenský fond za právnickou osobu pro tyto účely.\\
---\\
Možnost docílení pomocí označení svěřenského fondu jako právnické osoby v Občanském zákoníku a označení svěřenského fondu jako účastníka v Zákonu o zvláštních řízeních soudních.\\
---\\
Příkladem budiž Zákon o dani z přidané hodnoty.\\
---\\
\textit{"(2) Pro účely tohoto zákona se na svěřenský fond, podílový fond, podfond investičního fondu a organizační složku státu, která je účetní jednotkou, hledí jako na právnickou osobu."} (§4b, odstavec 2 ZDPH)

\end{frame}

% +=================+
% | Třináctý slide  |
% +=================+

\begin{frame}
	\frametitle{Zodpovězení otázek}
	\cevroLogo

\vspace*{-1cm}
\begin{table}[]
\begin{tabular}{|l|}
\hline
\begin{tabular}[c]{@{}l@{}}\textbf{§ 1478}\hspace{8.8cm}OZ\\ Za dědice nebo odkazovníka lze povolat i právnickou osobu,\\ která má teprve vzniknout. \\ Tato právnická osoba je způsobilým dědicem nebo odkazovníkem,\\ pokud vznikne do jednoho roku od smrti zůstavitele.\\ § \textbf{1478a}\\ Pro účely řízení o pozůstalosti se na svěřenský fond, hledí\\ jako na právnickou osobu.\end{tabular} \\ \hline
\end{tabular}
\end{table}

\begin{table}[]
\begin{tabular}{|l|}
\hline
\begin{tabular}[c]{@{}l@{}}\textbf{Účastníci při projednání pozůstalosti}\hspace{3.5cm}ZŘS\\ \textbf{§ 116}\\ \textbf{Svěřenský fond}\\ Svěřenský fond je účastníkem řízení, byl li zůstavitelem povolán\\ za dědice nebo v jeho prospěch nařídil odkaz. \\ Svěřenský fond zastupuje v průběhu pozůstalostního řízení\\ svěřenský správce.\end{tabular} \\ \hline
\end{tabular}
\end{table}

\end{frame}

% +================+
% | Čtrnáctý slide |
% +================+

\begin{frame}
	\frametitle{Zodpovězení otázek}
	\cevroLogo

\begin{table}[]
\begin{tabular}{|l|}
\hline
Otázka 2 Vedoucí práce                                                                                                                                                                                                                       \\ \hline
\begin{tabular}[c]{@{}l@{}}Lze změnit statut svěřenského fondu zřízeného mortis causa,\\ příp. za jakých podmínek? Kdo ho může změnit?\end{tabular} \\ \hline
Zřízení jako založení vs. vznik.\\
---\\
Zřízení (§ 1451) pro případ smrti - založen, zakladatel je živ.\\
---\\
Svěřenský fond neexistuje a proto je možné jej měnit stejně,\\ jako je možné měnit pořízení pro případ smrti.\\
---\\
Dědická dohoda - pouze po dohodě stran - závazek, například\\ novace.                                                                                                                                                                                                                              \\ \hline
\end{tabular}
\end{table}

\end{frame}

% +================+
% | Patnáctý slide |
% +================+

\begin{frame}
	\frametitle{Zodpovězení otázek}
	\cevroLogo

\begin{table}[]
\begin{tabular}{|l|}
\hline
Otázka 1 Oponent práce                                                                                                                                                                                                                       \\ \hline
\begin{tabular}[c]{@{}l@{}}Nechť autor v rámci ústní obhajoby shrne základní podobu\\ konceptu majetkové struktury svěřenského fondu.\end{tabular} \\ \hline
§495 OZ: \textit{"Souhrn všeho, co osobě patří, tvoří její majetek."}\\
---\\
Jedná se tedy o souhrn věcí, jejich příslušenství a součástí.\\
---\\
Majetek ve svěřenském fondu ale nikomu nepatří, není to\\ však ani věc ničí.\\
---\\
Jedná se o samostatně účelový majetek, který vlastní sám sebe.\\
---\\
Tento majetek spravuje svěřenský správce, který je rovněž\\ zapsaný ve všech příslušných registrech.\\
	\hline
\end{tabular}
\end{table}

\end{frame}

% +=================+
% | Šestnáctý slide |
% +=================+

\begin{frame}
	\frametitle{Zodpovězení otázek}
	\cevroLogo

\vspace*{-1cm}
\begin{table}[]
\begin{tabular}{|l|}
\hline
Otázka 2 Oponent práce                                                                                                                                                                                                                      \\ \hline
\begin{tabular}[c]{@{}l@{}}Autor by mohl dále shrnout některé dílčí odlišnosti québecké\\ a české úpravy svěřenských fondů a tím učinit závěr v dané otázce\\ na str. 20 až 23.\end{tabular} \\ \hline
V zásadě lze odlišnosti shrnout do několika následujících bodů.\\
---\\
Širší kategorizace účelů trustu v C.C.Q.\\
---\\
Private trust s určitým účelem, ne s určenou osobou, avšak\\ bez nutnosti být veřejně prospěšný, může být perpetuitou.\\
---\\
Absence statutu v C.C.Q., existence pouze \textit{trust instrumentu},\\ tedy samotné zakladatelské právní jednání.\\
---\\
Úprava kolize trustu s právními instituty\\ připodobnitelnými k právům nepominutelných dědiců.                                                                                                                                                                                                                        \\ \hline
\end{tabular}
\end{table}

\end{frame}

% +=================+
% | Sedmnáctý slide |
% +=================+

\begin{frame}

\begin{center}
\Large{Děkuji za pozornost	\smiley{}}
\end{center}

\end{frame}
	
\end{document}
